\begin{abstract}
Feed-forward 3D models such as VGGT and \pithree estimate cameras and geometry from a set of images in one forward pass, and their cost grows with the number of image tokens they process. Yet it is not clear how they use resolution: how their accuracy changes when the input is made larger or smaller, and which parts of the network actually need full resolution. We study these questions with frozen models and controlled probes. Each probe lowers the resolution at one stage of the network, such as the input pixels, the image encoder or global attention, for most of the views of a scene, and measures what is lost. We find that the two models respond to resolution in opposite ways: VGGT is accurate only near the input width it uses by default and degrades on both sides of it, while \pithree improves with resolution and then levels off. Resolution is not consumed at a single stage. For \pithree, lowering it in the encoder alone or in global attention alone costs most of what lowering it everywhere costs. Removing fine pixel detail while keeping the input size does not hurt \pithree, but it does hurt VGGT. In global attention, the views' queries need full resolution more than their keys. The models' errors at different input sizes are largely shared, so fusing predictions made at several sizes does not help. Finally, we show that 7-Scenes cannot resolve rotation differences below about 2\textdegree, which limits fine-grained comparisons on that benchmark.
\end{abstract}
